% Generated draft: requires mathematical and source review.
\documentclass[11pt]{article}
\usepackage[margin=25mm]{geometry}
\usepackage{fontspec}
\setmainfont{Latin Modern Roman}
\usepackage{amsmath,amssymb,mathtools,unicode-math}
\setmathfont{Latin Modern Math}
\usepackage{enumitem,needspace,xurl,hyperref}
\setlength{\parindent}{0pt}
\setlength{\parskip}{5pt}
\setlength{\emergencystretch}{3em}
\newcommand{\N}{\mathbb{N}}
\newcommand{\Z}{\mathbb{Z}}
\newcommand{\Q}{\mathbb{Q}}
\newcommand{\R}{\mathbb{R}}
\newcommand{\C}{\mathbb{C}}
\newcommand{\F}{\mathbb{F}}
\newcommand{\A}{\mathbb{A}}
\newcommand{\PP}{\mathbb{P}}
\newcommand{\E}{\mathbb{E}}
\newcommand{\eps}{\varepsilon}
\newcommand{\dd}{\,\mathrm{d}}
\begin{document}
\section*{UnsolvedMath ID: 30001131}
\textbf{Source problem code:} catalogue entry
\subsection{Definitions and mathematical statement}
For terminology, let \(k\) denote an unspecified base field. A homogeneous space in quotient notation is \(U=G/H\), where \(G\) is an algebraic group over \(k\), \(H\) is a closed subgroup, and the quotient exists as a variety. An equivariant compactification of \(U\) is a complete \(G\)-variety \(X\) together with a \(G\)-equivariant open immersion \(U\hookrightarrow X\) having dense image. These descriptions introduce notation; they do not impose a choice of field, group, subgroup, or regularity assumptions absent from the source.

In the local complete intersection sense, a closed subscheme \(Y\subseteq X\) is locally defined by a regular sequence: near each point of \(Y\), its ideal has generators \(f_1,\ldots,f_c\) such that multiplication by \(f_i\) is injective modulo \((f_1,\ldots,f_{i-1})\), and the generated ideal is proper. A global complete intersection instead requires suitable globally specified equations or sections. The source does not specify which meaning is intended.

\textbf{Source question.} Can the vanishing theory for complete intersections in toric varieties be generalized to complete intersections in algebraic homogeneous spaces and their equivariant compactifications?

\textbf{Unresolved formulation.} This is a request for a generalization of a theory, rather than a specified vanishing conjecture. The supplied record identifies neither the toric vanishing theorem to be generalized nor the cohomology theory, sheaves, degrees, positivity assumptions, or boundary conditions involved. For instance, an expression such as \(H^q(Y,\mathcal F)=0\) would require a specified sheaf \(\mathcal F\), degree range, and hypotheses; none are supplied. Consequently, no universal vanishing assertion, or precise class of allowable homogeneous spaces and compactifications, can be extracted without additional information.
\subsection{Short English statement}
The question asks whether vanishing results for complete intersections in toric varieties extend to algebraic homogeneous spaces and their equivariant compactifications. The intended vanishing results and hypotheses are unspecified, so the record supports a research question rather than a precise theorem. The supplied dated curation labels it partially solved; that assessment requires independent verification.
\subsection{Sources}
\begin{itemize}
\item Supplied local UnsolvedMath registry record, permanent ID \(30001131\), source problem code \texttt{catalogue entry}, titled \emph{Vanishing Theorems for Complete Intersections in Homogeneous Spaces}. The record gives proposed year 2009 but names no proposer.
\item \emph{Toric Geometry} (2009), Oberwolfach Reports, pp.\ 5--74, DOI 10.4171/owr/2009/01. Supplied source report and literature reference labeled ``original OWR report'': \url{https://doi.org/10.4171/owr/2009/01}. The background describes Oberwolfach Reports as documenting mathematical workshops and their open-problem discussions.
\item Supplied literature reference labeled ``Brion vanishing theorem'': \url{https://arxiv.org/abs/0812.2658}.
\item Literature assessment dated exactly \texttt{2026-08-21}: the supplied curation states that vanishing theorems for log homogeneous varieties and consequences for complete intersections have been proved, while the requested extension to complete intersections in all algebraic homogeneous spaces and compactifications is broader than the established results. Its supplied status is \texttt{partially\_solved}. This is supplied dated curation derived from the OWR material, not an independently verified assessment; the record expressly requests independent verification of time-sensitive status and cited claims.
\item The registry marks the statement as exact and reports a manual review finding the extracted statement self-contained, mathematically coherent, and free of LaTeX or report-layout artifacts. These are supplied verification claims, not verification performed for this draft, and do not resolve the missing mathematical specifications.
\end{itemize}
\end{document}
